\section*{الفصل \ref{ch:g7:plate-tectonics} --- تكتونية الصفائح: الأحجية المتحركة}

\begin{solution}{exo:g7:plate-tectonics:1}
\omterm{def:g7:plate-tectonics:continental-drift}{الانجراف القاري} هو الفكرة القائلة إن \omterm{def:g3:continents-and-oceans:continent}{القارات} تتحرك ببطء على
سطح الأرض، فتغيّرت مواقعها على مر العصور.
أدلة فيغنر: تطابق \omterm{def:g1:landscapes-around-us:water}{السواحل} (أمريكا الجنوبية وإفريقيا)؛
والأحافير نفسها على \omterm{def:g3:continents-and-oceans:continent}{قارات} تفصل بينها اليوم محيطات؛ \omterm{def:g3:stones-pebbles-sand:rock}{والصخور} نفسها \omterm{def:g4:world-physical-map:range}{والسلاسل
الجبلية} القديمة نفسها على جانبي \omterm{def:g2:water-on-earth:ocean}{محيط}؛ وآثار مناخات قديمة في غير
مكانها (الفحم في القارة القطبية الجنوبية، ورواسب جليدية في إفريقيا الحارة الجافة).
\end{solution}

\begin{solution}{exo:g7:plate-tectonics:2}
المتباعد: الصفيحتان تتباعدان. المتقارب: تقتربان إحداهما من الأخرى.
التحويلي: تنزلقان إحداهما بمحاذاة الأخرى، جنبًا إلى جنب.
\end{solution}

\begin{solution}{exo:g7:plate-tectonics:3}
مثلًا، حول \omterm{def:g2:water-on-earth:ocean}{المحيط} الهادئ كله (اليابان، وإندونيسيا، والأنديز)،
والحزام الممتد من \omterm{def:g2:water-on-earth:ocean}{البحر} المتوسط إلى الهيمالايا. ولا تكاد توجد
\omterm{def:g4:volcanoes-earthquakes-first-look:earthquake}{زلازل} كبيرة في وسط إفريقيا أو أستراليا (ولا في شرق أمريكا
الشمالية شرقي \omterm{def:g1:landscapes-around-us:landform}{جبال} روكي، ولا في البرازيل).
\end{solution}

\begin{solution}{exo:g7:plate-tectonics:4}
لا يستطيع حيوان يعيش في \omterm{def:g2:water-on-earth:salt-fresh}{الماء العذب} أن يعبر آلاف الكيلومترات من
\omterm{def:g2:water-on-earth:ocean}{المحيط} المالح. فوجوده على الجانبين يعني أن المكانين كانا ملتحمين يومًا
باليابسة، في \omterm{def:g3:continents-and-oceans:continent}{قارة} واحدة، حين كان \emph{الميزوسور} يعيش.
\end{solution}

\begin{solution}{exo:g7:plate-tectonics:5}
يتكون الفحم من نباتات غابات دافئة رطبة مستنقعية: فقد كان \omterm{def:g5:weather-and-climate:climate}{المناخ}
دافئًا رطبًا. ولا يمكن لمثل هذه الغابات أن تنمو قرب \omterm{def:g3:continents-and-oceans:poles}{القطب الجنوبي} اليوم، لذلك
لا بد أن القارة القطبية الجنوبية كانت يومًا أقرب بكثير إلى \omterm{def:g3:continents-and-oceans:equator}{خط الاستواء}: لقد تحركت.
\end{solution}

\begin{solution}{exo:g7:plate-tectonics:6}
ظهرة وسط الأطلسي: متباعد. الأنديز: متقارب (\omterm{def:g7:plate-tectonics:plate}{صفيحة} محيطية تغوص
تحت \omterm{def:g3:continents-and-oceans:continent}{قارة}). صدع سان أندرياس: تحويلي. الهيمالايا: متقارب (قارتان
تصطدمان). أخدود شرق إفريقيا: متباعد (\omterm{def:g3:continents-and-oceans:continent}{قارة} تنقسم).
\end{solution}

\begin{solution}{exo:g7:plate-tectonics:7}
$2.5 \times 10 = 25$ كيلومترًا في المليون سنة؛ وفي 60 مليون سنة:
$25 \times 60 = \qty{1500}{km}$.
\end{solution}

\begin{solution}{exo:g7:plate-tectonics:8}
$\qty{7}{m} = \qty{700}{cm}$، و$700 \div 9 \approx 78$: نحو
\qty{78}{cm} في السنة، أي نحو 30 ضعف المتوسط البالغ \qty{2.5}{cm} في السنة.
فالصفيحتان لا تنزلقان متباعدتين بانتظام عند السطح: بل يتراكم الإجهاد
سنوات ثم يتحرر في دفعات قصيرة، حين تنفتح الشقوق بعدة أمتار
دفعة واحدة. وعلى مدى فترات طويلة، تتعادل السنوات الهادئة والدفعات في متوسط بضعة
سنتيمترات في السنة.
\end{solution}

\begin{solution}{exo:g7:plate-tectonics:9}
لم يستطع أن يقول ما القوة التي تحرك \omterm{def:g3:continents-and-oceans:continent}{القارات}، وكانت الحركة التي
تصورها (\omterm{def:g3:continents-and-oceans:continent}{قارات} تشق طريقها عبر قاع \omterm{def:g2:water-on-earth:ocean}{المحيط}) مستحيلة.
وفي \omterm{def:g7:plate-tectonics:plate-tectonics}{تكتونية الصفائح}، لا تشق \omterm{def:g3:continents-and-oceans:continent}{القارات} طريقها عبر أي شيء: بل
تحملها، مع قاع \omterm{def:g2:water-on-earth:ocean}{المحيط} الذي يحيط بها، \omterm{def:g7:plate-tectonics:plate}{صفائح} تتحرك بوصفها قطعًا كاملة
من الغلاف الخارجي للأرض.
\end{solution}

\begin{solution}{exo:g7:plate-tectonics:10}
أعلى \omterm{def:g4:world-physical-map:altitude}{ارتفاع} عند الظهرة، وأعمق نقطة عند الخندق. تُصنع \omterm{def:g7:plate-tectonics:plate}{الصفيحة} الجديدة عند
الظهرة، حيث يصعد \omterm{def:g3:stones-pebbles-sand:rock}{صخر} حار من \omterm{def:g6:inside-the-earth:layers}{الوشاح}؛ وتعود \omterm{def:g7:plate-tectonics:plate}{الصفيحة} المحيطية إلى
\omterm{def:g6:inside-the-earth:layers}{الوشاح} عند الخندق، حيث تغوص تحت \omterm{def:g3:continents-and-oceans:continent}{القارة}.
\end{solution}

\begin{solution}{exo:g7:plate-tectonics:11}
$14.2 \div 2.2 \approx 6.45$: أسرع بنحو 6 مرات.
وبالكيلومترات في المليون سنة، السرعتان 142 و22. والمدتان اللازمتان لفتح
\qty{1000}{km}: $1000 \div 142 \approx 7$ ملايين سنة لظهرة شرق
\omterm{def:g2:water-on-earth:ocean}{المحيط} الهادئ؛ و$1000 \div 22 \approx 45$ مليون سنة لظهرة شمال
الأطلسي.
\end{solution}

\begin{solution}{exo:g7:plate-tectonics:12}
(أ) $6472 - 4578 = 1894$ \omterm{def:g4:volcanoes-earthquakes-first-look:earthquake}{زلزالًا}. (ب) $4578 \div 7031 \approx 0.651$،
أي 65\,\%. (ج) الخط المرسوم على \omterm{def:g2:drawing-to-map:map}{الخريطة} يحدد الخندق، حيث تبدأ \omterm{def:g7:plate-tectonics:plate}{الصفيحة} المحيطية
بالغوص. \omterm{def:g7:plate-tectonics:plate}{والصفيحة} لا تنزل عموديًا: بل تغوص
مائلة تحت \omterm{def:g7:plate-tectonics:plate}{الصفيحة} الأخرى، أعمق فأعمق كلما تقدمت تحتها.
وتقع \omterm{def:g4:volcanoes-earthquakes-first-look:earthquake}{الزلازل} على طول \omterm{def:g7:plate-tectonics:plate}{الصفيحة} الغائصة كلها، لذلك يقع كثير منها تحت
\omterm{def:g7:plate-tectonics:plate}{الصفيحة} العلوية، على الجانب نفسه من الخندق، على بعد عشرات أو مئات
الكيلومترات من الخط. \omterm{def:g7:plate-tectonics:plate-boundary}{فحد الصفائح} نطاق، لا خط رفيع.
\end{solution}

\begin{solution}{pb:g7:plate-tectonics:1}
\textbf{1.} \emph{الميزوسور} و\emph{الكينوغناثوس} و\emph{الغلوسوبتيريس}
(أما شريط \emph{الليستروسورس} فيصل إفريقيا بالهند والقارة القطبية الجنوبية).

\textbf{2.} كان يعيش في \omterm{def:g2:water-on-earth:salt-fresh}{الماء العذب}: فما كان ليستطيع أن يعبر محيطًا مالحًا
عرضه آلاف الكيلومترات.

\textbf{3.} يتطابق انتفاخ شمال شرق البرازيل مع تجويف
خليج غينيا، على \omterm{def:g1:landscapes-around-us:water}{الساحل} الغربي لإفريقيا.

\textbf{4.} \omterm{def:g7:plate-tectonics:plate-boundary}{حد متباعد}، هو ظهرة وسط الأطلسي: تتحرك \omterm{def:g7:plate-tectonics:plate}{صفيحة} أمريكا
الجنوبية نحو الغرب \omterm{def:g7:plate-tectonics:plate}{والصفيحة} الإفريقية نحو الشرق، مبتعدتين إحداهما عن الأخرى.

\textbf{5.} يقع الساحلان على جانبين متقابلين من \omterm{def:g6:latitude-longitude:longitude}{خط الزوال الأصلي}:
$43.2 + 14.4 = \ang{57.6}$.

\textbf{6.} $57.6 \times 102 = 5875.2$: نحو \qty{5880}{km}.

\textbf{7.} تتقارب \omterm{def:g6:latitude-longitude:longitude}{خطوط الزوال} كلما اتجهنا من
\omterm{def:g3:continents-and-oceans:equator}{خط الاستواء} نحو أحد \omterm{def:g3:continents-and-oceans:poles}{القطبين}. \omterm{def:g6:latitude-longitude:latitude}{ودائرة العرض} \ang{23} جنوبًا دائرة أصغر
من \omterm{def:g3:continents-and-oceans:equator}{خط الاستواء}، لذلك تكون كل درجة من درجاتها، وعددها 360، أقصر.

\textbf{8.} نحو نصف العرض: $5880 \div 2 = \qty{2940}{km}$.

\textbf{9.} $3.4 \times 10 = 34$ كيلومترًا في المليون سنة.

\textbf{10.} $5880 \div 34 \approx 173$: نحو 170 مليون سنة.

\textbf{11.} نعم: 170 مليون سنة تقع بين 100 و200 مليون سنة،
وهي تأتي بعد بدء تفكك \omterm{def:g7:plate-tectonics:supercontinent}{بانجيا}، قبل 225 إلى 200 مليون
سنة. فالحساب البسيط على \omterm{def:g2:drawing-to-map:map}{خريطة} يعطي رتبة
المقدار الصحيحة.

\textbf{12.} مثلًا: لم تكن السرعة واحدة طوال الوقت؛
وقد تمددت \omterm{def:g3:continents-and-oceans:continent}{القارات} وترققت قبل أن تنقسم، لذلك ليست \omterm{def:g1:landscapes-around-us:water}{سواحل} اليوم
هي الخط الذي انفصلت عنده تمامًا.
\end{solution}
